\documentclass[10pt,a4paper]{amsart}
\usepackage{amsmath,amssymb,amsthm,graphicx}
% Retained preamble for the Stacks Project chapter {CH} (shared project preamble).
% The chapter file pinned-original-source.tex only carries \input{{preamble}} on its first line,
% so no preamble text lives inside the chapter itself.  The Stacks Project uses one single shared
% preamble for every chapter: the snapshot stages it three times (sources/stacks-chow/preamble.tex,
% sources/stacks-etale-cohomology/preamble.tex and focused-seeds/AGP/sources/stacks-snapshot-20261011/
% preamble.tex) and those copies are identical in size.  The block below is copied mechanically,
% line for line, from sources/stacks-chow/preamble.tex lines 14-18 and 154-217: the Xy-pic setup and
% the native theorem declarations together with the mathematical macros.  Nothing was added, removed
% or redefined; the wrapper class and the amsmath/amssymb/amsthm/graphicx base packages are supplied
% by the assembler.
% For commutative diagrams we use Xy-pic
% Chapter: stacks-cohomology.tex -- Cohomology of Sheaves.
\usepackage[all]{xy}

% We use 2cell for 2-commutative diagrams.
\xyoption{2cell}

% Theorem environments.
%
\theoremstyle{plain}
\newtheorem{theorem}[subsection]{Theorem}
\newtheorem{proposition}[subsection]{Proposition}
\newtheorem{lemma}[subsection]{Lemma}

\theoremstyle{definition}
\newtheorem{definition}[subsection]{Definition}
\newtheorem{example}[subsection]{Example}
\newtheorem{exercise}[subsection]{Exercise}
\newtheorem{situation}[subsection]{Situation}

\theoremstyle{remark}
\newtheorem{remark}[subsection]{Remark}
\newtheorem{remarks}[subsection]{Remarks}

\numberwithin{equation}{subsection}

% Macros
%
\def\lim{\mathop{\mathrm{lim}}\nolimits}
\def\colim{\mathop{\mathrm{colim}}\nolimits}
\def\Spec{\mathop{\mathrm{Spec}}}
\def\Hom{\mathop{\mathrm{Hom}}\nolimits}
\def\Ext{\mathop{\mathrm{Ext}}\nolimits}
\def\SheafHom{\mathop{\mathcal{H}\!\mathit{om}}\nolimits}
\def\SheafExt{\mathop{\mathcal{E}\!\mathit{xt}}\nolimits}
\def\Sch{\mathit{Sch}}
\def\Mor{\mathop{\mathrm{Mor}}\nolimits}
\def\Ob{\mathop{\mathrm{Ob}}\nolimits}
\def\Sh{\mathop{\mathit{Sh}}\nolimits}
\def\NL{\mathop{N\!L}\nolimits}
\def\CH{\mathop{\mathrm{CH}}\nolimits}
\def\proetale{{pro\text{-}\acute{e}tale}}
\def\etale{{\acute{e}tale}}
\def\QCoh{\mathit{QCoh}}
\def\Ker{\mathop{\mathrm{Ker}}}
\def\Im{\mathop{\mathrm{Im}}}
\def\Coker{\mathop{\mathrm{Coker}}}
\def\Coim{\mathop{\mathrm{Coim}}}

% Boxtimes
%
\DeclareMathSymbol{\boxtimes}{\mathbin}{AMSa}{"02}

%
% Macros for moduli stacks/spaces
%
\def\QCohstack{\mathcal{QC}\!\mathit{oh}}
\def\Cohstack{\mathcal{C}\!\mathit{oh}}
\def\Spacesstack{\mathcal{S}\!\mathit{paces}}
\def\Quotfunctor{\mathrm{Quot}}
\def\Hilbfunctor{\mathrm{Hilb}}
\def\Curvesstack{\mathcal{C}\!\mathit{urves}}
\def\Polarizedstack{\mathcal{P}\!\mathit{olarized}}
\def\Complexesstack{\mathcal{C}\!\mathit{omplexes}}
% \Pic is the operator that assigns to X its picard group, usage \Pic(X)
% \Picardstack_{X/B} denotes the Picard stack of X over B
% \Picardfunctor_{X/B} denotes the Picard functor of X over B
\def\Pic{\mathop{\mathrm{Pic}}\nolimits}
\def\Picardstack{\mathcal{P}\!\mathit{ic}}
\def\Picardfunctor{\mathrm{Pic}}
\def\Deformationcategory{\mathcal{D}\!\mathit{ef}}

% --- appended 2026-10-12 (Chief): the project own definitions of the reference/slogan/history
% display environments, copied mechanically from the shared Stacks preamble lines 8-12
% (\usepackage{verbatim} plus the three \newenvironment lines, which the chapter preamble keeps
% verbatim from the project preamble). Nothing added, removed or redefined: these environments
% carry no mathematical content and the project itself typesets them as comments. Without these
% three lines a counted statement whose own text uses them fails with Environment slogan undefined.
\usepackage{verbatim}
\newenvironment{reference}{\comment}{\endcomment}
\newenvironment{slogan}{\comment}{\endcomment}
\newenvironment{history}{\comment}{\endcomment}
\usepackage{cleveref}
\crefname{originalreference}{original reference}{original references}
\Crefname{originalreference}{Original reference}{Original references}
\makeatletter
\newlabel{example-dual-derived}{{0FPC}{source}{}{sourcepointer-1}{}}
\newlabel{example-dual-derived@cref}{{[originalreference][0][]0FPC}{[1][0][]source}{}{}{}}
\makeatother
\crefname{originallemma}{original lemma}{original lemmas}
\Crefname{originallemma}{Original lemma}{Original lemmas}
\makeatletter
\newlabel{lemma-dual-perfect-complex}{{08DQ}{source}{}{sourcepointer-2}{}}
\newlabel{lemma-dual-perfect-complex@cref}{{[originallemma][0][]08DQ}{[1][0][]source}{}{}{}}
\makeatother
\crefname{originallemma}{original lemma}{original lemmas}
\Crefname{originallemma}{Original lemma}{Original lemmas}
\makeatletter
\newlabel{lemma-second-cup-equals-first}{{0FP2}{source}{}{sourcepointer-3}{}}
\newlabel{lemma-second-cup-equals-first@cref}{{[originallemma][0][]0FP2}{[1][0][]source}{}{}{}}
\makeatother
\crefname{originalequation}{original equation}{original equations}
\Crefname{originalequation}{Original equation}{Original equations}
\makeatletter
\newlabel{equation-eval}{{08I2}{source}{}{sourcepointer-4}{}}
\newlabel{equation-eval@cref}{{[originalequation][0][]08I2}{[1][0][]source}{}{}{}}
\makeatother
\crefname{originallemma}{original lemma}{original lemmas}
\Crefname{originallemma}{Original lemma}{Original lemmas}
\makeatletter
\newlabel{lemma-internal-hom-composition}{{0A8V}{source}{}{sourcepointer-5}{}}
\newlabel{lemma-internal-hom-composition@cref}{{[originallemma][0][]0A8V}{[1][0][]source}{}{}{}}
\makeatother
\title{The composition pairing between Ext and the dual equals the cup product under the identification of Hom with sections of a tensor product}
\author{The Stacks Project Authors}
\date{stacks-project:cohomology.tex:local-snapshot-20261011; GNU FDL 1.2 or later}
\begin{document}
\maketitle
\noindent\textit{Editorial note.} All delivered mathematical text is the author's own, copied verbatim from the pinned Stacks Project chapter file cohomology.tex; no statement or proof step was rewritten, reordered or removed.\par\medskip
\noindent\textit{Editorial note.} The statement contains a commutative diagram typeset with Xy-pic; the retained chapter preamble supplies the native Xy-pic setup, so it compiles exactly as in the source.\par\medskip
\noindent\textit{Editorial note.} The counted span cites nine cross-chapter labels (including an example and an equation), each recorded in assembly.external\_references with its canonical tag as a navigation pointer.\par\medskip
\noindent\textit{Editorial note.} The proof cites no bibliography item inside the counted span.\par\medskip
\noindent\textit{Allowed standard prerequisites (editorial).}\begin{itemize}
\item Comparison of the two cup products. Stacks Project, Cohomology of Sheaves, Lemma lemma-second-cup-equals-first (tag 0A8J). Identifies the cup product used in the statement with the tensor product of the representing maps.
\item Duality for perfect complexes: H\textasciicircum{}0(K tensor E\textasciicircum{}vee) = Hom(E, K). Stacks Project, Cohomology of Sheaves, Lemma lemma-dual-perfect-complex (tag 0A8C). Supplies the identification appearing as the left vertical arrow of the square.
\end{itemize}\medskip
\ifdefined\hypertarget\hypertarget{sourcepointer-1}{}\fi
\noindent\textit{Source reference.} Original reference 0FPC: sources/stacks-chow/tags.txt, source lines 19597--19597. Cross-chapter standard prerequisite cited by the author inside the counted span; it is used as a pointer and its own proof is not part of this item.\par\medskip
\ifdefined\hypertarget\hypertarget{sourcepointer-2}{}\fi
\noindent\textit{Source reference.} Original lemma 08DQ: sources/stacks-chow/tags.txt, source lines 10645--10645. Cross-chapter standard prerequisite cited by the author inside the counted span; it is used as a pointer and its own proof is not part of this item.\par\medskip
\ifdefined\hypertarget\hypertarget{sourcepointer-3}{}\fi
\noindent\textit{Source reference.} Original lemma 0FP2: sources/stacks-chow/tags.txt, source lines 19587--19587. Cross-chapter standard prerequisite cited by the author inside the counted span; it is used as a pointer and its own proof is not part of this item.\par\medskip
\ifdefined\hypertarget\hypertarget{sourcepointer-4}{}\fi
\noindent\textit{Source reference.} Original equation 08I2: sources/stacks-chow/tags.txt, source lines 10800--10800. Cross-chapter standard prerequisite cited by the author inside the counted span; it is used as a pointer and its own proof is not part of this item.\par\medskip
\ifdefined\hypertarget\hypertarget{sourcepointer-5}{}\fi
\noindent\textit{Source reference.} Original lemma 0A8V: sources/stacks-chow/tags.txt, source lines 12928--12928. Cross-chapter standard prerequisite cited by the author inside the counted span; it is used as a pointer and its own proof is not part of this item.\par\medskip
\begin{lemma}
\label{lemma-ext-composition-is-cup}
Let $(X, \mathcal{O}_X)$ be a ringed space. Let $K$ and $E$ be objects
of $D(\mathcal{O}_X)$ with $E$ perfect. The diagram
$$
\xymatrix{
H^0(X, K \otimes_{\mathcal{O}_X}^\mathbf{L} E^\vee) \times H^0(X, E)
\ar[r] \ar[d] &
H^0(X, K \otimes_{\mathcal{O}_X}^\mathbf{L} E^\vee
\otimes_{\mathcal{O}_X}^\mathbf{L} E) \ar[d] \\
\Hom_X(E, K) \times H^0(X, E) \ar[r] &
H^0(X, K)
}
$$
commutes where the top horizontal arrow is the cup product, the
right vertical arrow uses
$\epsilon : E^\vee \otimes_{\mathcal{O}_X}^\mathbf{L} E \to \mathcal{O}_X$
(Example \ref{example-dual-derived}), the left vertical arrow uses
Lemma \ref{lemma-dual-perfect-complex}, and the bottom horizontal
arrow is the obvious one.
\end{lemma}

\begin{proof}
We will abbreviate $\otimes = \otimes_{\mathcal{O}_X}^\mathbf{L}$
and $\mathcal{O} = \mathcal{O}_X$. We will identify $E$ and $K$
with $R\SheafHom(\mathcal{O}, E)$ and $R\SheafHom(\mathcal{O}, K)$
and we will identify $E^\vee$ with $R\SheafHom(E, \mathcal{O})$.

\medskip\noindent
Let $\xi \in H^0(X, K \otimes E^\vee)$ and $\eta \in H^0(X, E)$.
Denote $\tilde \xi : \mathcal{O} \to K \otimes E^\vee$ and
$\tilde \eta : \mathcal{O} \to E$ the corresponding maps in
$D(\mathcal{O})$. By Lemma \ref{lemma-second-cup-equals-first}
the cup product $\xi \cup \eta$ corresponds to
$\tilde \xi \otimes \tilde \eta : \mathcal{O} \to
K \otimes E^\vee \otimes E$.

\medskip\noindent
We claim the map $\xi' : E \to K$ corresponding to $\xi$ by
Lemma \ref{lemma-dual-perfect-complex} is the composition
$$
E = \mathcal{O} \otimes E
\xrightarrow{\tilde \xi \otimes 1_E}
K \otimes E^\vee \otimes E
\xrightarrow{1_K \otimes \epsilon}
K
$$
The construction in Lemma \ref{lemma-dual-perfect-complex}
uses the evaluation map (\ref{equation-eval}) which in turn
is constructed using the identification of $E$ with
$R\SheafHom(\mathcal{O}, E)$ and the composition
$\underline{\circ}$ constructed
in Lemma \ref{lemma-internal-hom-composition}.
Hence $\xi'$ is the composition
\begin{align*}
E = \mathcal{O} \otimes
R\SheafHom(\mathcal{O}, E)
& \xrightarrow{\tilde \xi \otimes 1}
R\SheafHom(\mathcal{O}, K) \otimes
R\SheafHom(E, \mathcal{O}) \otimes
R\SheafHom(\mathcal{O}, E) \\
& \xrightarrow{\underline{\circ} \otimes 1}
R\SheafHom(E, K) \otimes R\SheafHom(\mathcal{O}, E) \\
& \xrightarrow{\underline{\circ}}
R\SheafHom(\mathcal{O}, K) = K
\end{align*}
The claim follows immediately from this and the fact that
the composition $\underline{\circ}$ constructed in
Lemma \ref{lemma-internal-hom-composition} is associative
(insert future reference here) and the fact that $\epsilon$
is defined as the composition
$\underline{\circ} : E^\vee \otimes E \to \mathcal{O}$ in
Example \ref{example-dual-derived}.

\medskip\noindent
Using the results from the previous two paragraphs, we find
the statement of the lemma is that
$(1_K \otimes \epsilon) \circ (\tilde \xi \otimes \tilde \eta)$
is equal to
$(1_K \otimes \epsilon) \circ (\tilde \xi \otimes 1_E)
\circ (1_\mathcal{O} \otimes \tilde \eta)$
which is immediate.
\end{proof}

\end{document}
